% !PRINTED 198
\weberpage{198}
%<<CONT
\weq{6}{a_0 = 0, \quad a_1 = 1, \quad a_2 = 0, \quad a_3 = 0,}
so wird
\weq{7}{D = 0, \quad H = -x^2, \quad f = x^2 y, \quad Q = 2 x^3.}
Es würde also die linke Seite verschwinden, während die rechte
Seite nicht verschwindet, worin ein Widerspruch liegt.

Wir können aber endlich auch noch beweisen, dass auf der
linken Seite von (4) keine Potenz von 3 auftreten kann, die
nicht in allen Factoren $M$ enthalten ist, und sich also fortheben
lässt.

Wir haben im \S.\,2 den Satz bewiesen, dass das Product
zweier primitiver ganzer rationaler Functionen von beliebigen
Variablen wieder eine primitive Function ist. Dabei ist unter
einer primitiven Function eine solche verstanden, deren Coëfficienten
 ganze Zahlen ohne gemeinsamen Theiler sind. Es handelt
sich nun hier darum, nachzuweisen, dass eine Summe der Form
\[
Q^\tau \, \Sigma \, M D^\gamma H^\alpha f^\beta,
\]
worin die $M$ ganze Zahlen ohne gemeinsamen Theiler sind, nicht
imprimitiv werden, und speciell den Factor 3 erhalten kann,
wenn für $Q$, $D$, $H$, $f$ ihre Ausdrücke in den $a$, $x$, $y$ gesetzt
werden. Nach dem erwähnten Satze genügt es, da $Q$ eine primitive Function ist, dies nachzuweisen für die Form
\weq{8}{\Sigma \, M D^\gamma H^\alpha f^\beta.}
Hierin können wir überdies alle Glieder weglassen, deren $M$
den Factor 3 schon hat, und endlich können wir wieder nach
dem erwähnten Satze annehmen, dass der Ausdruck nicht durch
$D$ theilbar sei, dass er also ein Glied mit $\gamma = 0$ enthält.
Nehmen wir also an, es habe unter diesen Voraussetzungen der
entwickelte Ausdruck (8) den Theiler 3, und substituiren nun die
besonderen Werthe (6), (7), d.\,h. setzen wir $D = 0$, $H = -x^2$,
$f = x^2 y$, so reducirt sich der Ausdruck auf das einzige Glied,
in dem $\gamma = 0$, $2\alpha = \lambda$ ist,
\[
(-1)^\alpha M x^{2(\alpha + \beta)} y^\beta
\]
und es würde also folgen, dass dies $M$ den Theiler 3 haben
müsste, was gegen die Voraussetzung ist. Damit ist also bewiesen:

\gesperrt{Jede ganzzahlige Covariante der cubischen Form
kann ganz und rational und mit ganzzahligen Coëffi-}
%CONT>>
